\begin{prob}
    For each of the following functions, express its rate of growth using
    $\Theta$-notation, and prove your answer by finding constants which satisfy
    the definition of $\Theta$-notation. Make sure to show your work by writing
    out the chain of inequalities to prove each bound, like we did in lecture.
    Note: please do \emph{not} use limits to prove your answer (though you can
    use them to check your work).

    \begin{subprobset}
        \begin{subprob}
            $f(n) = 3n^2 - 8n + 4$

            \begin{soln}
                $f(n) = \Theta(n^2)$.

                We start with an upper bound:
                \begin{align*}
                    f(n)
                    &= 3n^2 - 8n + 4 \\
                    &\leq
                        3n^2 + 4
                        \\
                    &\leq
                        3n^2 + 4n^2 \qquad (n \geq 1)
                        \\
                    &= 7n^2
                \end{align*}

                Then for a lower bound:

                \begin{align*}
                    f(n) &= 3n^2 - 8n + 4\\
                         &\geq
                            3n^2 - 8n
                            \\
                        \intertext{Our problem now is the $-8n$. Remember that
                            we can do anything that makes the quantity
                            \emph{smaller} when trying to find a lower bound,
                            so we can't simply drop the $-8n$ since that would
                            actually make the quantity bigger. Instead, we recognize intuitively
                            that $n^2$ will eventually be larger than $-8n$, so we can
                            ``break off'' a piece of the $3n^2$ and use it to ``balance'' the
                            $-8n$.}
                        &=
                            2n^2 + n^2 - 8n
                            \\
                        &=
                            2n^2 + (n^2 - 8n)
                            \\
                        \intertext{
                            Again, our intuition is that $n^2 - 8n$ will eventually be
                            positive when $n$ is large enough, and so we'll be able to drop it to get a lower
                            bound. When is this? We solve $n^2 - 8n \geq 0$ for $n$,
                            finding that $n^2 \geq 8n \implies n \geq 8$. So when $n \geq 8$,
                            the stuff in parens is positive and can be dropped:
                        }
                        &\geq
                            2n^2 \qquad (n \geq 8)
                \end{align*}

                Therefore, $2n^2 \leq f(n) \leq 7n^2$ provided that $n \geq 8$.
            \end{soln}

        \end{subprob}

        \begin{subprob}
            $
            \displaystyle
            f(n) =
            \frac{n^2 + 2n - 5}{n - 10}
            $

            \begin{soln}
                $f(n) = \Theta(n)$.

                We begin with an upper bound.
                \begin{align*}
                    f(n)
                    &=
                        \frac{n^2 + 2n - 5}{n - 10}
                        \\
                    &\leq
                        \frac{n^2 + 2n}{n - 10}
                        \\
                    &=
                        \frac{n(n + 2)}{n - 10}
                        \\
                        \intertext{Now, $(n+2)/(n-10) \leq 2$ for all $n \geq 22$, hence:}
                    &\leq
                        2 n, \qquad (n \geq 22)
                        \\
                \end{align*}
                Now for a lower bound:
                \begin{align*}
                    f(n)
                    &=
                        \frac{n^2 + 2n - 5}{n - 10}
                        \\
                    &\geq
                        \frac{n^2 - 5}{n - 10}
                        \\
                    \intertext{To get a lower bound, we make this
                    quantity smaller. We can do so by making the
                    denominator bigger by dropping the $-10$:}
                    &\geq
                        \frac{n^2 - 5}{n}
                        \\
                    &\geq
                    \frac{\frac{1}{2}n^2  + \frac{1}{2}n^2 - 5}{n}
                        \\
                        \intertext{Now, $\frac{1}{2}n^2 -5 \geq 0$ for all $n \geq 4$. Hence:}
                    &\geq
                \frac{\frac{1}{2}n^2  + 0}{n}, \qquad (n \geq 4)
                        \\
                    &= n/2
                \end{align*}
                So in total, $\frac{n}{2} \leq f(n) \leq 2n$ for all $n \geq 22$.
            \end{soln}
            
        \end{subprob}

        \begin{subprob}
            $
            \displaystyle
            f(n) = 
            \begin{cases}
                n^2, & n \text{ is odd},\\
                4n^2, & n \text{ is even}
            \end{cases}
            $

            \begin{soln}
                $f(n) = \Theta(n^2)$.
                
                Our goal is to find a function that upper bounds $f$ both when $n$ is even and when it is odd, and a separate
                function that lower bounds $f$ both when $n$ is even and when it is odd.

                In this case, $f(n) \leq 4n^2$ for all $n \geq 1$, whether $n$ is even or odd.
                Likewise, $f(n) \geq n^2$ for all $n \geq 1$, even or odd. Hence:
                $n^2 \leq f(n) \leq 4n^2$ for all $n \geq 1$.
                
                You might be tempted to say that the first part, $n^2$, is $\Theta(n^2)$, and 
                the second part, $4n^2$, is also $\Theta(n^2)$, so the whole function is $\Theta(n^2)$.
                This \emph{is} a valid argument, but we haven't proven that it's true -- though you can
                prove it is valid using the techniques that
                we learned in lecture, but it's a little too ``high level'' for this problem. Here, we
                want to go all the way back to the definition.
            \end{soln}
        \end{subprob}

        \begin{subprob}
            State the growth of the function below using $\Theta$ notation in as
            simplest of terms possible, and prove your answer by finding constants
            which satisfy the definition of $\Theta$ notation.

            E.g., if $f(n)$ were $3n^2 + 5$, we would write $f(n) = \Theta(n^2)$ and not $\Theta(3n^2)$.

            \[
                f(n) = \frac{n^3 - n^2 + n + 1000}{(n-1)(n+2)}
            \]
            
            \begin{soln}
                $\Theta(n)$.

                Recall that we write $f(n) = \Theta(g(n))$ if there exist positive constants $c_1, c_2, N$
                such that
                \[
                    c_1 g(n) \leq f(n) \leq c_2 g(n)
                \]
                for all $n \geq N$.

                In this case, we will prove the inequality for $g(n) = n$. There are infinitely-many choices
                of $c_1, c_2, $ and $N$ that will satisfy the inequality; we will prove such one:

                \underline{Upper Bound} :  $f(n) \leq c_2 \cdot g(n)$ \\

                \begin{align*}
                \frac{n^3-n^2 + n + 1000}{(n-1)(n+2)}
                & \le \frac{1002n^3}{n^2 + n - 2}\\
                & \le \frac{1002n^3}{n^2 + (n - 2)} (n \ge ?) \\
                & \le \frac{1002n^3}{n^2} (n \ge 2) \\
                & \le 1002n \\
                \end{align*}


                \underline{Lower Bound} : $c_1 \cdot g(n) \leq f(n)$ \\

                \begin{align*}
                \frac{n^3-n^2 + n + 1000}{(n-1)(n+2)}
                & \ge \frac{n^3-n^2}{n^2 + n}\\
                & \ge \frac{0.5n^3 + (0.5n^3 -n^2)}{n^2 + n} (n \ge ?) \\
                & \ge \frac{0.5n^3}{n^2 + n} (n \ge 2) \\
                & \ge \frac{0.5n^3}{n^2 + n^2} \\
                & \ge \frac{1}{4}n \\
                \end{align*}

                We have $n \ge 2 $ for both the upper and lower bound. Hence, $N=2$. \\
                Therefore $f(n) = \theta(n)$ with constants $N=2, c_1 = 0.25, c_2 = 1002$
            \end{soln}
        \end{subprob}

    \end{subprobset}
\end{prob}
